\documentclass[11pt]{article}
\usepackage[a4paper,margin=1in]{geometry}
\usepackage{amsmath,amssymb}
\usepackage[hidelinks]{hyperref}
\setlength{\parindent}{0pt}
\setlength{\parskip}{0.65em}
\emergencystretch=3em

\newcommand{\Lk}{\Lambda_k}

\title{Literature Status for Higher-Rank Numerical Ranges\\
\large arXiv:math/0511278 answered by arXiv:0706.1536}
\author{Run \texttt{fa\_banach\_001}, lane 4}
\date{August 11, 2026}

\begin{document}
\maketitle

\textbf{Status.} \texttt{literature\_already\_answered}. This note records
an exact later theorem, not a new proof of the run.

\section*{Source questions}
For $A\in M_n(\mathbb C)$ and $1\leq k\leq n$, Choi--Kribs--Życzkowski
define
\[
 \Lk(A)=\{\mu\in\mathbb C:PAP=\mu P\text{ for some rank-$k$
 orthogonal projection }P\}.
\]
Conjecture~2.8 of arXiv:math/0511278 asks whether, for a normal matrix with
eigenvalues $\lambda_1,\ldots,\lambda_n$ counted with multiplicity,
\[
 \Lk(A)=
 \bigcap_{1\leq j_1<\cdots<j_{n-k+1}\leq n}
 \operatorname{conv}\{\lambda_{j_1},\ldots,\lambda_{j_{n-k+1}}\}.
 \tag{CKŻ}
\]
Problem~2.9 asks more generally whether $\Lk(A)$ is convex whenever it is
nonempty. The paper identifies the first unresolved normal case as
$(n,k)=(5,2)$ and asks whether all interior points of the pentagonal region
for the cyclic five-shift belong to $\Lambda_2$.

\section*{Complete later answer}
Li--Sze, arXiv:0706.1536, Theorem~2.2, prove for every complex matrix $A$
that
\[
 \Lk(A)=\bigcap_{\xi\in[0,2\pi)}
 \left\{\mu\in\mathbb C:
 e^{i\xi}\mu+e^{-i\xi}\overline\mu
 \leq \lambda_k\!\left(e^{i\xi}A+e^{-i\xi}A^*\right)
 \right\},
 \tag{LS}
\]
where $\lambda_k(H)$ denotes the $k$th largest eigenvalue of a Hermitian
matrix $H$. Every set in the intersection is a closed half-plane.
Corollary~2.3 therefore gives
\[
 \boxed{\Lk(A)\text{ is convex for every }A\in M_n(\mathbb C).}
\]
This answers Problem~2.9 affirmatively, including the empty-set convention
and hence slightly strengthening the source's phrasing.

For normal $A$, Li--Sze Corollary~2.4 reduces formula~\textup{(LS)} to
formula~\textup{(CKŻ)} verbatim. Thus Conjecture~2.8 is also completely
affirmative. In particular, the cyclic five-shift's $\Lambda_2$ is the full
intersection of the five four-eigenvalue convex hulls: every interior point
of the inner pentagon is included.

\section*{Identification boundary}
The intermediate paper arXiv:quant-ph/0608244 verified the conjecture for a
wide class of unitary and normal matrices. Li--Sze's Theorem~2.2 and
Corollaries~2.3--2.4 are the decisive general result: they cover arbitrary
matrices for convexity and every finite normal matrix for the exact spectral
formula. No additional hypothesis or parameter translation is needed.

\begin{thebibliography}{9}
\bibitem{CKZ}
M.-D. Choi, D. W. Kribs, and K. Życzkowski,
\emph{Higher-Rank Numerical Ranges and Compression Problems},
arXiv:math/0511278; Linear Algebra Appl. 418 (2006), 828--839.

\bibitem{LS}
C.-K. Li and N.-S. Sze,
\emph{Canonical forms, higher rank numerical range, convexity, totally
isotropic subspace, matrix equations},
arXiv:0706.1536; Proc. Amer. Math. Soc. 136 (2008), 3013--3023,
Theorem~2.2 and Corollaries~2.3--2.4.
\end{thebibliography}

\end{document}
